% Insertar bajo el encabezado 2.6 Implementación del documento principal.
% Requiere en el preámbulo: \usepackage{fancyvrb}

La solución se organiza en una plataforma web de gestión clínica, una API
de servicios, una base de datos relacional y un componente móvil orientado
a la consulta del odontólogo. Esta separación permite concentrar las
operaciones administrativas en la web y delimitar la experiencia móvil a
los pacientes asignados y sus estudios radiográficos. CERPAX participa como
proveedor externo mediante un portal de entrega con enlace temporal.

\subsubsection{Plataforma web y gestión clínica}

La interfaz web se implementó con Angular y TypeScript. Los componentes de
Angular Material se utilizaron para construir formularios, tablas,
diálogos de confirmación y elementos de navegación. La aplicación consume
los servicios de la API mediante solicitudes HTTP y presenta las acciones
disponibles según los permisos de la sesión.

Se implementaron operaciones de registro, consulta y actualización de
pacientes y odontólogos, así como la gestión de horarios y citas. En las
citas se valida que el intervalo solicitado corresponda al horario del
profesional y se rechazan las superposiciones. Las consultas incorporan
paginación y búsqueda para evitar la transferencia de todos los registros
en una sola respuesta.

Los pacientes y los odontólogos comparten la tabla \texttt{usuarios}, pero
su función se distingue mediante las asociaciones del sistema. La tabla
\texttt{pacientes\_odontologos} relaciona al paciente con el profesional y
la clínica correspondiente. Esta asignación se utiliza para delimitar los
registros que puede consultar cada odontólogo. El paciente se representa
como un registro clínico y no como un actor con acceso al panel.

\subsubsection{Autenticación y control de acceso}

El backend se desarrolló con Node.js y Express. El inicio de sesión de la
plataforma web se realiza mediante \texttt{/api/login}; la contraseña se
verifica con bcrypt y, si las credenciales son válidas, se emite un token
JWT con vigencia de ocho horas. Las solicitudes protegidas incluyen el
token en el encabezado de autorización.

El middleware de autenticación comprueba la firma, la vigencia del token,
el estado del usuario y la versión de su sesión. Posteriormente se verifica
el permiso requerido y la pertenencia a la clínica seleccionada. De esta
manera, ocultar una opción en la interfaz no constituye el único control:
las restricciones también se aplican en la API y en las consultas de datos.

El administrador dispone de funciones de administración de cuentas y
permisos. El odontólogo consulta la información correspondiente a sus
pacientes asignados. El acceso de recepción se limita a datos
administrativos, citas y asignaciones, excluyendo los campos clínicos y
las operaciones de administración de cuentas.

El siguiente fragmento muestra la comprobación de permisos que utiliza
el backend. La función \texttt{can} comprueba el permiso de la operación
o el rol de administrador; el middleware rechaza la solicitud con el
código HTTP 403 cuando la comprobación falla. La excepción
\texttt{bypass} corresponde al acceso local de desarrollo, condicionado
por la configuración del servidor, y no habilita el acceso local en producción.

\noindent\textbf{Fragmento 1. Control de permisos por operación.}
\begin{Verbatim}[fontsize=\footnotesize,frame=single,numbers=left]
export const can = (user, key, action) => !!user && (
  user.bypass || user.roles.includes('Super Admin') ||
  user.permissions.includes(`${key}.${action}`)
);
export function requirePermission(key, action) {
  return (req, res, next) => {
    try {
      if (!can(req.user, key, action)) {
        fail('No tienes permiso para realizar esta acción', 403);
      }
      next();
    } catch (e) { next(e); }
  };
}
\end{Verbatim}
\noindent\textit{Fuente: código del proyecto,
\texttt{backend/src/auth.js}; formato ajustado para su lectura.}

\subsubsection{Solicitudes y entrega de radiografías}

El odontólogo registra una solicitud radiográfica indicando el paciente,
el tipo de estudio y las indicaciones. La API comprueba la asignación
antes de crear el registro. Las solicitudes utilizan los estados
\textit{Pendiente}, \textit{Completada} y \textit{Cancelada}.

Para la entrega de CERPAX se implementó la generación de un enlace y un
código QR con vigencia configurable entre una y setenta y dos horas. El
token se genera mediante bytes aleatorios y la base de datos conserva su
hash, junto con las fechas de vencimiento, revocación y utilización. Al
generar un nuevo acceso se revocan los enlaces anteriores que todavía no
han sido utilizados.

El proveedor externo no necesita iniciar sesión en el panel. El portal
valida el token y solicita una URL firmada para cargar el archivo en
Supabase Storage. Una vez transferido el archivo, la API verifica sus
metadatos y la correspondencia con el acceso temporal. La confirmación
registra la radiografía, completa la solicitud y marca el acceso como
utilizado dentro de una transacción.

Se admiten archivos PDF, PNG, JPEG y WebP con un tamaño máximo de
\mbox{50 MiB}. Los objetos se almacenan en un bucket privado y sus
metadatos se conservan en PostgreSQL. La visualización o descarga desde el
panel requiere autorización; conocer el identificador de una solicitud
no habilita por sí mismo el acceso al archivo.

El segundo fragmento muestra la validación del acceso temporal. La consulta
compara el hash del token y exige que el enlace no haya vencido, sido revocado
o utilizado, y que la solicitud permanezca pendiente. Si no se encuentra un
acceso válido, la API devuelve el código HTTP 410. El parámetro
\texttt{lock} permite bloquear las filas durante las operaciones que
modifican el estado de la solicitud dentro de una transacción.

\noindent\textbf{Fragmento 2. Validación del enlace temporal de CERPAX.}
\begin{Verbatim}[fontsize=\footnotesize,frame=single,numbers=left]
async function access(client, token, lock = false) {
  const {rows} = await client.query(
    `SELECT a.id AS acceso_id,a.vence_en,a.storage_path,s.*
     FROM accesos_radiografias a JOIN solicitudes_radiografias s
       ON s.id=a.solicitud_id
     WHERE a.token_hash=$1
       AND a.vence_en>now()
       AND a.revocado_en IS NULL
       AND a.usado_en IS NULL
       AND s.estado='Pendiente'
     ${lock ? 'FOR UPDATE OF s,a' : ''}`,
    [tokenHash(token)]
  );
  if (!rows[0]) fail(
    'El enlace venció, fue revocado o ya fue utilizado',
    410
  );
  return rows[0];
}
\end{Verbatim}
\noindent\textit{Fuente: código del proyecto,
\texttt{backend/src/radiography.js}; formato ajustado para su lectura.}

\subsubsection{Componente móvil para el odontólogo}

La aplicación móvil DentalCare utiliza Flutter y Dart para la interfaz
Android. Dentro de la solución, su alcance se centra en el inicio y cierre
de sesión del odontólogo, la consulta de pacientes asignados, la consulta
del estado de las solicitudes y la visualización de las radiografías
disponibles. Esta delimitación permite utilizar el teléfono como medio de
consulta clínica sin trasladar todas las operaciones del panel web.

El diseño de integración establece la API de la plataforma como punto de
acceso a los datos y como responsable de verificar la identidad, los
permisos y las asignaciones. La administración de cuentas, roles y
permisos, así como la gestión de odontólogos, horarios y citas, se concentra
en la plataforma web. El paciente no inicia sesión en la aplicación móvil
y CERPAX mantiene el acceso mediante el enlace temporal de entrega.

\subsubsection{Decisiones técnicas y resolución de dificultades}

Una dificultad técnica del flujo radiográfico consiste en mantener la
coherencia entre el archivo almacenado y la solicitud clínica. Para
atenderla, se separaron la preparación de la carga, la transferencia del
archivo y la confirmación en la API. La solicitud se completa después de
verificar el objeto cargado; los cambios relacionados en PostgreSQL se
ejecutan en una transacción. La restricción de unicidad de
\texttt{solicitud\_id} en \texttt{radiografias} impide registrar más de
una radiografía para la misma solicitud.

Otra dificultad consiste en restringir el acceso cuando distintos
profesionales utilizan el mismo sistema. Se resolvió incorporando filtros
por clínica y por asignación de pacientes en las consultas del backend,
además de las comprobaciones de permisos por operación. Las solicitudes
fuera del alcance autorizado se rechazan o se excluyen de los resultados.

Para delimitar el alcance funcional se retiraron del catálogo y de la
navegación los módulos de contabilidad, laboratorio y departamentos, junto
con los roles de contador y laboratorista. Las tablas heredadas se
conservaron para evitar la pérdida de información y mantener la
compatibilidad con registros existentes. Esta decisión distingue la
retirada de una función visible de la eliminación física de sus datos.

Los cambios del esquema se gestionan mediante migraciones SQL versionadas.
Las consultas emplean parámetros para los valores recibidos y los nombres
de recursos se resuelven desde el catálogo permitido. Asimismo, los errores
de validación se presentan mediante respuestas estructuradas, mientras que
los errores internos evitan revelar credenciales o detalles de conexión.

\subsubsection{Verificación de la implementación}

La verificación del backend incluyó pruebas de autenticación, validación
de datos, gestión clínica, migraciones, restricciones de recepción y flujo
radiográfico. Se comprobaron, entre otros escenarios, el acceso de un
odontólogo a pacientes no asignados, los enlaces vencidos o revocados,
la carga única y las restricciones entre clínicas. En la ejecución local
documentada se aprobaron veinticinco pruebas y se completó la compilación
de producción de Angular. Estos resultados corresponden al backend y a la
aplicación web verificados, y no constituyen una medición de rendimiento
ni una prueba de ejecución de la aplicación móvil.

El código completo, las migraciones y las pruebas se mantienen en el
repositorio del proyecto. Los diagramas de arquitectura, datos y casos
de uso complementan la descripción de los componentes y sus relaciones.
